\section{Metric Definitions}
\label{app:metric_defs}

This appendix gives the full definitions of every metric recorded by the
evaluation pipeline of Section~\ref{subsec:metrics}. The main text explains the
measures the argument turns on; the tables below are the look-up reference for
the complete set.

\paragraph{Stability-gap metrics.} Computed from the dense per-step evaluation.
Write $a_j^{\text{pre}}$ for the pre-task accuracy on task $j$ and $a_j(t)$ for
its accuracy at step $t$ of the new task.

\begin{table*}[htbp]
    \centering
    \small
    \begin{tabular}{@{}p{2.3cm} p{7.9cm} p{4.4cm}@{}}
        \hline
        Metric                   & Definition                                                                                                              & Role                                                 \\
        \hline
        \texttt{gap\_depth}      & $\max_j \left( a_j^{\text{pre}} - \min_t a_j(t) \right)$ over the full record window                                    & Primary metric --- worst single point of instability \\
        \texttt{gap\_area}       & $\sum_j \int_0^{T} \max\!\left(0,\, a_j^{\text{pre}} - a_j(t)\right) dt$, trapezoidal over the full new-task record (to task end $T$); accuracy at or above baseline contributes $0$ & Depth $\times$ duration of the below-baseline shortfall (units: accuracy $\cdot$ steps) \\
        \texttt{max\_drop}       & \texttt{gap\_depth} restricted to the first $200$ steps                                                                 & Fixed-window metric, kept for back-comparability     \\
        \texttt{recovery\_steps} & Given a record has dipped below the threshold, the first step --- counted from the transition --- at which \emph{every} past task is simultaneously at $\geq 90\%$ of its pre-task baseline & None if no sub-threshold dip occurs, or if recovery is not reached within the record \\
        \hline
    \end{tabular}
    \caption{Stability-gap metrics, recorded from dense per-step evaluation. \texttt{gap\_depth} is the headline measure used in every contrast.}
    \label{tab:gap_metrics}
\end{table*}

\paragraph{Continual-learning metrics.} Following
\citet{delange2023continualevaluationlifelonglearning}, computed at task
boundaries and over the full record. Let $A(E_j, f_n)$ be the accuracy on the
test set of task $j$ evaluated at the model $f_n$ after training stage $n$, and
let $N$ be the number of tasks. The full task-boundary accuracy matrix
$R[i,j] = A(E_j, f_{T_i})$ is logged for every run, so any downstream metric can
be reconstructed.

\begin{table*}[htbp]
    \centering
    \small
    \begin{tabular}{@{}p{1.8cm} p{8.5cm} p{4.3cm}@{}}
        \hline
        Metric          & Definition                                                                                                      & Role                                              \\
        \hline
        ACC             & $\frac{1}{N} \sum_j A(E_j, f_{T_{N-1}})$                                                                        & Average task accuracy at the final model          \\
        FORG            & $\frac{1}{N-1} \sum_{i<N-1} \left[ A(E_i, f_{T_i}) - A(E_i, f_{T_{N-1}}) \right]$                               & Average drop from right-after-learning to final   \\
        min-ACC         & $\frac{1}{N-1} \sum_{i<N-1} \min\{ A(E_i, f_n) : n > T_i \}$                                                    & Worst-case retention since learning               \\
        $\mathrm{WF}_w$ & Mean over tasks of the largest accuracy drop in any window of $w$ consecutive evaluations ($w \in \{10, 100\}$) & Worst-case stability over short / medium horizons \\
        $\mathrm{WP}_w$ & Symmetric counterpart of $\mathrm{WF}_w$: the largest accuracy gain ($w \in \{10, 100\}$)                       & Plasticity / learning velocity                    \\
        WC-ACC          & $\frac{1}{N} A(E_{N-1}, f_{T_{N-1}}) + \left(1 - \frac{1}{N}\right)\text{min-ACC}$                              & Worst-case trade-off --- a lower bound on ACC     \\
        \hline
    \end{tabular}
    \caption{Continual-learning metrics of \citet{delange2023continualevaluationlifelonglearning}, recorded at task boundaries and over the full record.}
    \label{tab:cl_metrics}
\end{table*}

\paragraph{Gradient-fidelity diagnostics.} Recorded at the dense cadence for the
decomposition conditions only. The tracker compares the replay-buffer gradient
$g_{\text{replay}}$ against the true past-task gradient $g_{\text{true}}$,
computed at every step by a separate forward and backward pass over the entire
past-task training set (no sub-sampling, so $g_{\text{true}}$ is the
deterministic empirical past-task gradient at the current parameters).

\begin{table*}[htbp]
    \centering
    \small
    \begin{tabular}{@{}p{4.2cm} p{3.0cm} p{7.2cm}@{}}
        \hline
        Diagnostic                                    & Logged name                     & Role                                                                                               \\
        \hline
        $\cos(g_{\text{replay}}, g_{\text{true}})$    & \texttt{true\_grad\_cosine}     & Buffer-side directional fidelity --- the estimator-bias contributor of Section~\ref{subsec:decomp} \\
        $\|g_{\text{replay}}\| / \|g_{\text{true}}\|$ & \texttt{true\_grad\_mag\_ratio} & Buffer-side magnitude fidelity                                                                     \\
        $\|g_{\text{new}}\| / \|g_{\text{replay}}\|$  & \texttt{grad\_ratio}            & Update-side magnitude asymmetry --- large at step $1$ when $\|g_{\text{replay}}\| \approx 0$       \\
        \hline
    \end{tabular}
    \caption{Gradient-fidelity diagnostics, recorded over the new task's first ${\sim}500$ steps. Each is logged as a per-step series together with its mean (and minimum, for the cosine). The first two carry the \texttt{true\_grad} prefix because they are measured against $g_{\text{true}}$; \texttt{grad\_ratio} compares $g_{\text{replay}}$ to the current-task gradient $g_{\text{new}}$ and its inverse drives the adaptive curriculum of Section~\ref{subsec:curriculum}.}
    \label{tab:grad_metrics}
\end{table*}
